\chapter{Последнее рассеяние и бариогенез}

\section{Смотреть вдаль — значит смотреть в прошлое}

На разных расстояниях мы видим Вселенную в разные моменты времени. На малых расстояниях время, на которое мы заглядываем в прошлое, не очень существенно. Вселенная год назад или сто лет назад не слишком отличалась от Вселенной сейчас. Но на достаточно больших расстояниях мы видим Вселенную более молодой. Более молодая Вселенная плотнее, а более плотная — горячее.

На некотором достаточно большом расстоянии Вселенная была примерно такой же горячей, как поверхность Солнца. Нам она так не выглядит, потому что фотоны по пути к нам испытали красное смещение. Коэффициент красного смещения порядка тысячи означает наблюдаемую температуру, меньшую во столько же раз:
\[
T_{\mathrm{obs}} \sim \frac{T_{\mathrm{emit}}}{1+z},
\qquad
1+z \sim 10^3.
\]\footnote{Эта запись представляет собой сжатую редакторскую передачу словесного рассуждения о красном смещении из 00:02:20--00:02:39.}

Именно там Вселенная стала непрозрачной. Она состояла главным образом из водорода, с небольшой примесью гелия. При температуре выше нескольких тысяч градусов водород ионизован. Поэтому по ту сторону этой поверхности находятся не нейтральные атомы, а ионы: главным образом свободно движущиеся протоны и электроны. Фотоны там тоже более энергичны, по существу это рентгеновские фотоны. Если атом случайно образуется, столкновения или фотон быстро снова его ионизуют.

Ионизованный газ является электрическим проводником, а проводник не даёт излучению проходить беспрепятственно. Заряды реагируют на электрические поля, колеблются взад и вперёд и переизлучают его. Поэтому дальше этого расстояния Вселенная для нас непрозрачна. Поверхность последнего рассеяния подобна видимой поверхности Солнца: это не материальная стена, а место, где среда довольно резко переходит от непрозрачности к прозрачности.

Большинство частиц там — фотоны: их примерно в \(10^8\) раз больше. То же верно и сегодня, хотя сегодня энергия фотонов настолько мала, что они не составляют важной части современной плотности энергии. Когда температура падает ниже температуры ионизации водорода, Вселенная становится прозрачной, и мы можем видеть сквозь неё.

На этом прямое наблюдение заканчивается. Следующий вопрос состоит в том, насколько дальше в прошлое нас ещё позволяют продвинуться известные атомная, ядерная физика и физика элементарных частиц.

\section{Ещё дальше в прошлое}

Сразу по ту сторону поверхности последнего рассеяния, помимо фотонов, основными частицами являются протоны и электроны, а также некоторое количество ядер гелия и немного дейтронов. Антипротонов практически нет, и позитронов тоже практически нет. Но дальше в прошлое температура в конце концов становится достаточно высокой, чтобы два фотона могли столкнуться и породить электрон и позитрон:
\[
\gamma + \gamma \leftrightarrow e^- + e^+.
\]

Основной порог определяется тем, что фотоны должны по меньшей мере обеспечить энергию покоя пары:
\[
E \ge 2m_e c^2.
\]
Ниже этого порога обычного света недостаточно. Выше него могут рождаться электрон--позитронные пары, и обратная реакция тоже может происходить. В горячей порции вещества при такой температуре происходили бы реакции между частицами, играющие ту же роль, что химические реакции в обычном веществе.

Непосредственно выше порога фотонов всё ещё больше, чем электронов и позитронов. Масса покоя электрона энергетически дорога. Но когда температура становится намного выше порога, энергия покоя электрона оказывается менее важной, чем кинетическая энергия, и численные плотности становятся сравнимыми. Тогда электроны, позитроны и фотоны имеют примерно одинаковые численные плотности и сталкиваются с бешеной частотой.

\begin{figure}[ht]
\centering
\includegraphics[width=0.62\textwidth]{lecture_06_figure_02.png}
\caption{Схема порога рождения пар для сравнения масштабов масс электрона и протона.}
\lectureframe{00:18:47}
\end{figure}

Следующий вопрос — что происходит с протонами. Они уже присутствуют, и к этому времени их численная плотность очень велика, но тепловой энергии ещё недостаточно, чтобы в изобилии рождать протон--антипротонные пары. Причина проста:
\[
m_p \approx 2000\,m_e,
\qquad
E \ge 2m_p c^2
\]
— это гораздо более высокий порог.\editorialnote{На кадре нижняя подпись массы менее чёткая, чем верхняя \(2m_e c^2\), но в устном объяснении сравнение явно проводится между электронным и протонным порогами.}

Это уже намного горячее, чем внутри звезды, где античастицы не присутствуют в сколько-нибудь заметном количестве.

Но и это лишь предварительное описание. К тому времени, когда температура становится достаточно высокой, чтобы протоны и антипротоны свободно рождались, она уже достаточно высока и для того, чтобы разбивать протоны на кварки. Поэтому эту по-настоящему горячую среду лучше описывать как смесь электронов, позитронов, кварков, антикварков, фотонов, нейтрино, глюонов и других видов частиц. Простая схематическая реакция с доски такова:
\[
e^- + e^+ \to \gamma \to q + \bar q.
\]

\begin{classroomqa}
\lecturetimestamp{00:20:05}
\audiencequestion{Это мезоны или прекварки?}
\lecturerresponse{Мезоны состоят из кварков. Как только температура становится достаточно высокой, чтобы разрушить протон, она разрушает и мезоны.}
\end{classroomqa}

До сих пор рассуждение шло назад во времени. Теперь развернём его и спросим, что происходит, когда горячая симметричная смесь остывает.

\section{Электрическая нейтральность и горячая симметричная Вселенная}

При достаточно высокой температуре во Вселенной имеются колоссальные концентрации частиц и античастиц. Электроны и позитроны появляются парами, кварки и антикварки появляются парами, и естественно ожидать почти равных количеств каждой частицы и её античастицы.

Понятие электрической нейтральности требует некоторой осторожности. Сегодня Вселенная электрически нейтральна. Полезный мнемонический приём — представить не нашу реальную Вселенную, а замкнутую и ограниченную Вселенную, например трёхсферу. Сначала рассмотрим в качестве модели двумерную сферу. Силовые линии электрического поля не могут начинаться или заканчиваться нигде, кроме как на заряде. Положительный заряд испускает силовые линии; отрицательный их поглощает. На замкнутой сфере такие линии не могут просто уйти на бесконечность. Они снова сходятся на противоположной стороне. Поэтому в качестве мнемонического образа можно сказать, что замкнутая и ограниченная Вселенная не может иметь ненулевой суммарный заряд. Дело не в том, является ли наша реальная Вселенная в буквальном смысле замкнутой. Суть в том, что на больших масштабах положительные и отрицательные заряды, по-видимому, компенсируют друг друга чрезвычайно точно.

Электрическая нейтральность \emph{не} означает, что число протонов равно числу антипротонов. Она означает, что полный положительный заряд равен полному отрицательному заряду. Протоны, позитроны и все прочие положительно заряженные частицы вместе уравновешивают электроны и все прочие отрицательно заряженные частицы вместе.

Настолько далеко назад во времени, насколько можно доверять известной физике, и, возможно, немного дальше, Вселенная была достаточно горячей, чтобы в ней в изобилии существовали и частицы, и античастицы. Любому ускорительному эксперименту, какой только можно назвать, на каком-то этапе соответствовал аналог в ранней Вселенной. Почти симметричное начало — естественная гипотеза: равные числа электронов и позитронов, равные числа кварков и антикварков и, если протоны и антипротоны использовать в качестве представителей барионов и антибарионов,
\[
N_p = N_{\bar p}.
\]

Эта гипотеза привлекательна, потому что законы физики почти симметричны по отношению к частицам и античастицам. Позитрон ведёт себя почти так же, как электрон, не совсем, но очень близко. Именно эта почти симметрия и есть проблема.

\section{Прокручивая фильм вперёд}

Предположим, что Вселенная начинается в горячем состоянии с равными или почти равными количествами частиц и античастиц, а затем остывает.

Пока температура достаточно высока, протоны и антипротоны могут рождаться в энергичных столкновениях, и они могут аннигилировать обратно в фотоны. Но в конце концов температура опускается ниже протонного порога. Тогда протон--антипротонные пары ещё могут аннигилировать в фотоны, но после некоторого расширения эти фотоны теряют достаточно энергии и уже не могут снова превращаться в протон--антипротонные пары. Это улица с односторонним движением. То же самое позже происходит с электронами и позитронами, при температуре гораздо ниже протонного порога, но всё ещё намного выше звёздных температур.

Поскольку расширение происходит медленно по сравнению со скоростями микроскопических реакций, во время этого процесса Вселенная большую часть времени остаётся близкой к тепловому равновесию. Именно это здесь и имеется в виду под адиабатическим расширением. В этих условиях точнее говорить, что энтропия меняется очень мало. Для горячей смеси частиц энтропия по существу определяется числом частиц, с точностью до множителей порядка единицы. Так что протоны, антипротоны, электроны и позитроны, которые аннигилируют, приблизительно заменяются фотонами.

Если бы числа протонов и антипротонов были в точности равны, они просто полностью аннигилировали бы. То же самое произошло бы с электронами и позитронами. Остались бы фотоны. Но если бы существовал хотя бы крошечный избыток протонов над антипротонами, все образовавшиеся пары аннигилировали бы, а неспаренный избыток остался бы. Конкретный пример проясняет мысль: если бы было сто протонов и пятьдесят антипротонов, то эти пятьдесят антипротонов нашли бы пятьдесят протонов, аннигилировали бы с ними, и осталось бы пятьдесят протонов. Тогда нейтральность заряда требует такого же рода оставшегося избытка электронов над позитронами.

Это позволяет восстановить раннюю асимметрию по нынешнему числу фотонов. Сегодня число фотонов примерно от \(10^9\) до \(10^{10}\) раз больше числа протонов или электронов:
\[
\frac{N_\gamma}{N_p} \sim 10^9\text{--}10^{10}.
\]
Следовательно, в раннем горячем супе на каждый миллиард протон--антипротонных пар приходился лишь примерно один лишний протон. Именно из-за этого крошечного дисбаланса атомы до сих пор существуют.

\begin{classroomqa}
\lecturetimestamp{00:42:25}
\audiencequestion{Почему не рассматривать этот избыток статистически, как нечто, возникающее при остывании с очень малой вероятностью?}
\lecturerresponse{Чтобы получить лишний протон из первоначально симметричного состояния, пришлось бы создать протон, не создав одновременно антипротон. Вместе с ним пришлось бы создать и отрицательный заряд, скажем электрон. В известной стандартной модели такого низкоэнергетического процесса нет. При тех проверенных энергиях барионное число сохраняется, так что обычное остывание не порождает чистого избытка.}
\end{classroomqa}

Итак, электрический заряд не является препятствием. Схематический пример таков:
\[
\gamma + \gamma \to p + e^-.
\]\footnote{Эта запись представляет собой редакторскую расшифровку словесного примера из 00:44:16--00:46:35. Ее цель — выделить мысль о том, что электрический заряд мог бы сохраняться при нарушении барионного числа.}
В целом это было бы электрически нейтрально. Но здесь нарушается другое эмпирическое правило. В проверенной стандартной модели нельзя создать протон, не создав одновременно антипротон. Для наших целей важной сохраняющейся величиной является барионное число,
\[
B=(N_p+N_n)-(N_{\bar p}+N_{\bar n}).
\]\footnote{Эта компактная формула формализует словесное определение, данное в 00:45:18--00:45:40.}
При известных энергиях эта величина сохраняется, так что асимметрия не может быть создана на поздней, хорошо проверенной части истории остывания.

\section{Бариогенез}

Именно этим и занимается бариогенез. Предположим, что Вселенная действительно началась с равных количеств барионов и антибарионов. Что должны были бы допускать законы физики, чтобы сместить это симметричное состояние к наблюдаемому крошечному избытку?

На доске это было сведено к двум символическим выражениям,
\[
N_p=N_{\bar p},
\qquad
N_p-N_{\bar p},
\]
а затем — к списку условий, связывающих первое со вторым.

\begin{figure}[ht]
\centering
\includegraphics[width=0.72\textwidth]{lecture_06_figure_04.png}
\caption{Условия порождения барионной асимметрии из первоначально симметричного состояния.}
\lectureframe{01:22:26}
\end{figure}

\editorialnote{На этом кадре слова в нумерованном списке читаются лишь частично. Приведенные ниже три пункта следуют сопутствующему устному обсуждению.}

Первое условие — нарушение барионного числа. Если барионное число является точным законом сохранения, то Вселенная, которая начинается с нулевого суммарного барионного числа, такой и остается.

Второе условие — наличие предпочтения в пользу частиц по сравнению с античастицами. Одного лишь нарушения барионного числа недостаточно. Если природа допускает процесс, порождающий протон, и соответствующий ему процесс, порождающий антипротон, с равной вероятностью, то средний дисбаланс все равно равен нулю. Нужно, чтобы скорости этих процессов различались.

Обычные статистические флуктуации тоже слишком малы. Если полное число частиц, включая фотоны, составляет около \(10^{90}\), то флуктуация порядка \(\sqrt{N}\) дает лишь около \(10^{45}\), что и близко не достаточно, чтобы объяснить примерно \(10^{80}\) протонов во Вселенной.

Третье условие состоит в том, что на каком-то этапе Вселенная должна была быть достаточно далека от равновесия. Здесь в рассуждении используется теорема CPT. Если горячий газ вечно находился в равновесии, то фильм, пущенный назад, выглядит так же, как фильм, пущенный вперед. В такой ситуации симметрия между частицами и античастицами эффективно восстанавливается благодаря точной симметрии, объединяющей зарядовое сопряжение, пространственную четность и обращение времени. Чтобы реальное предпочтение закрепилось, фактическая история должна на микроскопических временных масштабах выглядеть по-разному при движении вперед и назад во времени. На достаточно раннем этапе расширение Вселенной было достаточно быстрым, чтобы это обеспечить.

Итак, три условия таковы:
\begin{enumerate}
\item нарушение барионного числа;
\item нарушение симметрии между частицами и античастицами, точнее, CP-нарушение;
\item стадия вне теплового равновесия.
\end{enumerate}

\begin{classroomqa}
\lecturetimestamp{01:09:08}
\audiencequestion{Если соотнести это с CPT, то условие 2 — действительно CP, а условие 3 — действительно вопрос о направлении времени?}
\lecturerresponse{Да. Условие 2 на самом деле следует формулировать как CP-нарушение. Условие 3 — это то обстоятельство, что фактическая история была вне равновесия, а не отдельный закон сохранения и не то же самое, что нарушение барионного числа.}
\end{classroomqa}

При наличии этих трех условий даже Вселенную, которая начинается с равных чисел протонов и антипротонов, можно склонить в ту или иную сторону. Неясным остается то, действительно ли в природе нарушается барионное число и если да, то как.

\section{Что известно, а что нет}

Нерешенной частью этой истории остается нарушение барионного числа. Мы никогда непосредственно не наблюдали в природе реакцию, изменяющую
\[
(N_p+N_n)-(N_{\bar p}+N_{\bar n}).
\]
Это может означать лишь то, что соответствующие реакции требуют гораздо более высоких температур, чем все, что когда-либо непосредственно проверялось в лаборатории, или же что при низкой энергии они фантастически маловероятны. Но на данный момент именно здесь и остается знак вопроса.

Вот почему в этом обсуждении так важен распад протона. Если барионное число не является точным, то протон в принципе должен быть нестабилен. Но он поразительно стабилен. Приводимая здесь нижняя граница имеет порядок
\[
\tau_p \gtrsim 10^{34}\ \text{years}.
\]
Она получена в гигантских экспериментах с очень большими объемами воды, оснащенными световыми детекторами и расположенными глубоко под землей. Это не опровергает нарушение барионного числа. Это означает лишь, что любое такое нарушение при низкой энергии должно быть чрезвычайно сильно подавлено.

В обсуждение входит и еще одно замечание из известной физики. При энергиях от \(100\) до \(250\,\mathrm{GeV}\) слабые взаимодействия становятся гораздо более похожими на электромагнитные. Это может иметь значение в детальных моделях, но не является решающим шагом в данном рассуждении.

Темная материя не слишком меняет основную логику. В большинстве теорий она взаимодействует с обычным веществом настолько слабо, что для данного вопроса по существу остается сторонним наблюдателем.

\begin{classroomqa}
\lecturetimestamp{01:24:49}
\audiencequestion{Как на все это влияет гравитация?}
\lecturerresponse{Главным образом двояко. Во-первых, без космического расширения не было бы стадии вне равновесия. Во-вторых, испарение черных дыр дает серьезное основание подозревать, что барионное число не может вечно быть точным законом сохранения. Кроме этого, обычные гравитационные взаимодействия между частицами слишком слабы, чтобы здесь существенно влиять.}
\end{classroomqa}

Затем возникает уточняющий вопрос об электрическом заряде. Ставит ли тот же аргумент с черными дырами под угрозу и сохранение заряда? Ответ: нет. Заряженная черная дыра несет линии электрического потока, выходящие наружу, и по ним сохраняется информация о заряде. Такая черная дыра не испаряется до конца; вместо этого она приближается к экстремальному остатку. Следовательно, аргумент с черной дырой ставит под угрозу барионное число, а не электрический заряд.

Другой вопрос состоит в том, не следует ли просто считать очень плотный ранний кварковый суп черной дырой. Ответ: нет. Речь идет не о статическом плотном сгустке. Вселенная расширялась. История с черной дырой — это мысленный эксперимент о судьбе барионного числа, а не буквальное описание ранней Вселенной как запертой области.

\begin{classroomqa}
\lecturetimestamp{01:31:14}
\audiencequestion{Понятно ли, как модифицировать стандартную модель, чтобы включить такие реакции?}
\lecturerresponse{Трудность как раз противоположная. Слишком легко добавить чрезмерное нарушение барионного числа и тем самым заставить протон распадаться чрезмерно быстро. Настоящая проблема состоит в том, чтобы сохранить протон столь же стабильным, как он есть, и при этом все же допустить нужные процессы ранней Вселенной.}
\end{classroomqa}

\section{Энтропия и поздние вопросы о нейтронах}

Теперь вернемся к учету. В адиабатически расширяющейся Вселенной энтропия приблизительно сохраняется. Для горячего супа из элементарных частиц энтропия по существу совпадает с числом частиц с точностью до множителей порядка единицы. Поэтому в этом обсуждении не требуется никакого огромного позднего рождения фотонов из ничего. Аннигиляции в основном лишь перераспределяют содержимое этого супа в излучение.

В более старой формулировке проблемы все было поставлено с ног на голову. Вместо вопроса о том, почему по сравнению с фотонами уцелело так мало протонов, задавали обратный: почему фотонов так много по сравнению с протонами? Но более точный вопрос состоит не в том, почему Вселенная на позднем этапе породила так много энтропии. Он состоит в том, почему асимметрия между частицами и античастицами с самого начала была столь мала: лишь примерно один лишний протон на миллиард протон--антипротонных пар.

Свободный нейтрон действительно распадается, но этот процесс \emph{не} нарушает барионное число:
\[
n \to p + e^- + \bar\nu.
\]
Нейтрон несет то же барионное число, что и протон.

Для свободного нейтрона доступная энергия распада очень мала, поскольку масса нейтрона лишь немного превышает массу конечного состояния. Отчасти поэтому по меркам физики элементарных частиц этот распад происходит так медленно. Когда нейтрон связан в ядре, энергия связи может полностью устранить эту небольшую доступную энергию, и тогда связанный нейтрон может быть стабильным.

\begin{classroomqa}
\lecturetimestamp{01:43:03}
\audiencequestion{Нейтроны просто исчезают из Вселенной, или они также образуются?}
\lecturerresponse{Первый ответ состоит в том, что свободные нейтроны теряются вследствие распада. Но этим дело не исчерпывается. Существуют реальные источники образования нейтронов, включая ядерные реакции в звёздах и образование нейтронных звёзд. Так что общее число нейтронов не просто однонаправленно убывает.}
\end{classroomqa}
